\section{Agentic RL 基础设施的更高难度}

\subsection{从静态验证到动态环境}

当优化目标从解决 benchmark 问题转向解决交互式任务时，RL 技术栈必须随之改变。Reasoning RL 中 rollout 通常可以作为基本自包含的轨迹处理，配合相对干净的评估器。而在 Agentic RL 中，策略（policy）被嵌入一个更大的工具系统中：工具服务器、浏览器、终端、搜索引擎、模拟器、执行沙箱、API 层、记忆系统和编排框架。

\begin{warningbox}{环境不再是静态验证器}
在 agentic RL 中，\textbf{环境本身成为训练系统的一部分}。这创造了一个新的系统需求：训练和推理必须更干净地解耦。否则，rollout 吞吐量将崩溃。
\end{warningbox}

\subsection{训练-推理解耦的必要性}

作者给出了一个生动的例子：考虑一个编码 agent 必须针对实时测试工具执行生成的代码——推理端因等待执行反馈而停滞，训练端因缺少完成的轨迹而挨饿，整个流水线的 GPU 利用率远低于经典 reasoning RL 的预期。工具延迟、部分可观测性和有状态环境进一步放大了这些低效。

\subsection{环境成为一等研究工件}

\begin{importantbox}{从数据多样性到环境质量}
\textit{In the SFT era, we obsessed over data diversity. In the agent era, we should obsess over environment quality.}（在 SFT 时代，我们痴迷于数据多样性。在 agent 时代，我们应该痴迷于环境质量。）环境质量包括：稳定性、真实性、覆盖度、难度、状态多样性、反馈丰富度、抗利用性，以及 rollout 生成的可扩展性。
\end{importantbox}

环境构建已经从一个副项目变成了一个真实的创业类别。如果 agent 被训练在类生产环境中运行，那么环境就是核心能力栈的一部分。

\subsection{本章小结}

Agentic RL 的基础设施比 reasoning RL 困难得多：策略嵌入复杂工具系统，训练-推理必须解耦以避免吞吐量崩溃，环境本身成为一等研究工件。行业需要从``痴迷数据多样性''转向``痴迷环境质量''。


